% Part: counterfactuals
% Chapter: introduction
% Section: paradoxes-material

\documentclass[../../../include/open-logic-section]{subfiles}

\begin{document}

\olfileid{cnt}{int}{par}

\olsection{مفارقات الشرط المادي}

كان سي.~آي. لويس من أوائل من انتقدوا استعمال~$!A \lif !B$ وسيلةً
لتمثيل العبارات الإنجليزية من قبيل «إذا \dots، فإن \dots». وكان لويس
ينتقد استعمال الشرط المادي في كتاب وايتهد وراسل
\emph{مبادئ الرياضيات}، حيث كانا يقرآن~$\lif$ بمعنى «يستلزم».
وقد اعترض لويس بحق على أنه إذا كان~$\lif$ يعني «يستلزم»، فإن كل
قضية كاذبة~$p$ تستلزم أن~$p$ تستلزم~$q$، لأن
$p \lif (p \lif q)$ صادقة إذا كانت~$p$ كاذبة؛ كما أن كل قضية صادقة~$q$
تستلزم أن~$p$ تستلزم~$q$، لأن~$q \lif (p \lif q)$ صادقة إذا كانت~$q$
صادقة.

يعرف المناطقة بالطبع أن \emph{الاستلزام}، أي الاستلزام المنطقي، ليس
رابطًا، بل علاقة بين !!{formula} أو العبارات. ولذلك ينبغي ألّا نقرأ
$\lif$ بمعنى «يستلزم»، اتقاءً للالتباس.\footnote{ما زالت قراءة
  «$\lif$» بمعنى «يستلزم» واسعة الانتشار بين الرياضيين وعلماء الحاسوب،
  مع أن الفلاسفة يحاولون تفادي الالتباسات التي أبرزها لويس بقراءته
  بمعنى «فقط إذا».} وما دمنا لا نفعل ذلك، فإن موضع قلق لويس بعينه لا
ينشأ أصلًا: فـ$p$ لا «تستلزم»~$q$ حتى لو اعتبرنا~$p$ رمزًا لجملة
إنجليزية كاذبة بالفعل. ولكي نقرر هل~$p \Entails q$، علينا أن ننظر في
\emph{جميع} !!{valuation}، ويظل~$p \Entails/ q$ حتى حين نستعمل~$p$
لتمثيل جملة تصادف أنها كاذبة.

ومع ذلك، يبقى شيء مستغرب في عبارات «إذا \dots، فإن \dots» من قبيل
مثال لويس:
\begin{quote}
إذا كان القمر مصنوعًا من الجبن الأخضر، فإن~$2+2=4$.
\end{quote}
وكذلك في الاستدلالين:
\begin{quote}
  القمر ليس مصنوعًا من الجبن الأخضر. ولذلك، إذا كان القمر مصنوعًا
  من الجبن الأخضر، فإن~$2+2=4$.

  $2+2 = 4$. ولذلك، إذا كان القمر مصنوعًا
  من الجبن الأخضر، فإن~$2+2=4$.
\end{quote}
ومع هذا، لو كانت «إذا \dots، فإن \dots» هي~$\lif$ فحسب، لكانت الجملة
صادقة بلا إشكال، ولكان الاستدلالان صالحين بلا إشكال.

ويتعلق مثال آخر بالقضية الصادقة منطقيًا
$(!A \lif !B) \lor (!B \lif !A)$. وهذا يوحي بأنه إذا أخذنا عشوائيًا
جملتين خبريتين~$S$ و$T$ من صحيفة، فينبغي أن تكون الجملة «إذا كانت~$S$
فإن~$T$، أو إذا كانت~$T$ فإن~$S$» صادقة.

\end{document}
